\section{Conclusion}
\label{sec:conclusion}

We presented GraphManager, a retrieval framework that leverages an AST-derived graph world model to provide cost-aware file localization for software issue resolution. The framework supports seven retrieval methods ranging from a deterministic graph scorer with zero LLM runtime cost to LLM-guided interactive agents, all evaluated under a dual-track protocol that separates strict commit-fidelity from same-snapshot amortization.

Our experiments on SWE-bench issues across multiple Python repositories show that graph-guided retrieval achieves competitive file-level F1 compared to strong RAG baselines while requiring substantially fewer total tokens. The graph index setup costs approximately 3$\times$ fewer embedding tokens than dense code chunk indexing, and this advantage compounds in same-snapshot settings where the graph serves multiple retrieval tasks. Paired per-issue analysis with bootstrap confidence intervals confirms that the quality differences are statistically meaningful on completed experiment cells.

These results are preliminary: the evaluation matrix is partially incomplete, sample sizes are small, and we evaluate retrieval quality without downstream patch generation. Completing the full seven-repository matrix, expanding to larger issue samples, and adding execution-based validation are immediate priorities.

Looking forward, two extensions are particularly promising. First, incremental graph updates---recomputing only the affected subgraph when files change---could further reduce setup cost for evolving repositories. Second, extending the graph construction pipeline to additional languages via tree-sitter's multi-language grammars would broaden the applicability of graph-guided retrieval to the multi-language codebases represented in benchmarks such as SWE-PolyBench \citep{swepolybench2025}.
